\documentclass[12pt]{article}
\usepackage[utf8]{inputenc}
\usepackage{graphicx}
\usepackage{amsmath}
\usepackage{amssymb}
\usepackage{mathrsfs}
\usepackage{pgfornament}
\usepackage[many]{tcolorbox}
\usepackage[margin = 0.5in]{geometry}
\usepackage{collectbox}
\usepackage{mdframed}
\usepackage{xcolor}
\usepackage{mathtools}
\DeclarePairedDelimiter\ceil{\lceil}{\rceil}
\DeclarePairedDelimiter\floor{\lfloor}{\rfloor}
\newcommand\norm[1]{\left\lVert#1\right\rVert}
\newcommand\textbox[1]{%
  \parbox{.333\textwidth}{#1}%
}

\linespread{2}

\makeatletter
\newcommand{\mybox}{%
    \collectbox{%
        \setlength{\fboxsep}{2pt}%
        \fbox{\BOXCONTENT}%
    }%
}
\makeatother

\usepackage{listings}
\usepackage{color}

\definecolor{dkgreen}{rgb}{0,0.6,0}
\definecolor{gray}{rgb}{0.5,0.5,0.5}
\definecolor{mauve}{rgb}{0.58,0,0.82}

\lstset{frame=tb,
  language=R,
  aboveskip=3mm,
  belowskip=3mm,
  showstringspaces=false,
  columns=flexible,
  basicstyle={\small\ttfamily},
  numbers=none,
  numberstyle=\tiny\color{gray},
  keywordstyle=\color{blue},
  commentstyle=\color{dkgreen},
  stringstyle=\color{mauve},
  breaklines=true,
  breakatwhitespace=true,
  tabsize=3
}
\title{MAT 523 HW 7}
\begin{document}

\hfill March 29th, 2019
\begin{center}
\begin{large}
\textbf{MAT 524: Functions of a Real Variable II \\
\vspace{-5mm}
Homework VII \\}
\end{large}
\vspace{-2mm}
Nolan R. H. Gagnon \\
\end{center}
\vspace{-0.5cm}
\noindent \textbf{{\huge [}{\Large [}} Problem Three \textbf{{\Large ]}{\huge ]}}{\Large \textbf{:}}

\vspace{3mm}

\noindent\mybox{\colorbox{yellow}{\textbf{Problem Statement}}}\hspace{3mm}\hrulefill

\noindent Using Bessel's inequality, prove that for every measurable subset $A \subseteq [0,1]$,
$$
\lim\limits_{n\rightarrow\infty}\int_A \cos(2\pi n x) dx = \lim\limits_{n\rightarrow\infty} \sin(2\pi n x)dx = 0.
$$

\noindent\mybox{\colorbox{yellow}{\textbf{Solution}}}\hspace{3mm}\hrulefill

\noindent We begin by noting that the function $y = \frac{\chi_A}{\sqrt{2}}$ is in $L^2([0,1])$ and the set $$\bigcup\limits_{n=1}^{\infty}\left\lbrace \sqrt{2}\cos(2\pi n x)\right\rbrace$$ is orthonormal in $L^2([0,1])$. For $n\in\mathbb{N}$, the Fourier coefficients of $y$ are 
\begin{align}
\hat{y}(n) &= \left<y,\sqrt{2}\cos(2\pi n x)\right> \\
&= \int\limits_0^1 \frac{\chi_A}{\sqrt{2}} \sqrt{2}\cos(2\pi n x) dx \\
&= \int\limits_{A} \cos(2\pi n x) dx.
\end{align}\
By Bessel's inequality,
$$\sum\limits_{n=1}^{\infty} \left| \int_A \cos(2\pi n x) \right|^2 \leq \norm{y}^2.$$ But
\begin{align}
    \norm{y}^2 &= \left(\left[\int\limits_0^1 \left| \frac{\chi_A}{\sqrt{2}}\right|^2\right]^{1/2}\right)^2 \\
    &= \int_A \frac{1}{\sqrt{2}} \\
    &= \frac{m(A)}{\sqrt{2}} \\
    &< \frac{1}{\sqrt{2}}.
\end{align}
Thus, the sum $\sum\limits_{n=1}^{\infty} \left| \int_A \cos(2\pi n x) \right|^2$ converges, so we must have
$$\lim\limits_{n\rightarrow\infty} \left| \int_A \cos(2\pi n x) \right|^2 = 0.$$ Since the squared absolute value function is continuous, we can move the limit inside, obtaining
$$\left|\lim\limits_{n\rightarrow\infty} \int_A \cos(2\pi n x) \right|^2=0,$$ i.e.,
$$\lim\limits_{n\rightarrow\infty} \int_A \cos(2\pi n x) = 0.$$

Now, since $$\bigcup\limits_{n=1}^{\infty}\left\lbrace \sqrt{2}\sin(2\pi n x)\right\rbrace$$ is also orthonormal in $L^2([0,1])$, we can apply the same argument as above to conclude that $$\lim\limits_{n\rightarrow\infty} \int_A \sin(2\pi n x) = 0.$$


\hfill$\blacksquare$

\vfill

\begin{center}
\pgfornament[width = 50mm, color = black, symmetry = h]{83}\\
\vspace{2mm}
\vspace{-0.65cm}
\hrulefill \\
\vspace{-0.95cm}
\hrulefill
\end{center}

\end{document}
